\section{Results}

% ---------------------------------------------------------------------------
%  This file is a worked pattern, not filler. Replace the prose; keep the
%  shapes. Every construct below is verified to export cleanly to Word.
% ---------------------------------------------------------------------------

\subsection{First result}

\todo{prose}

% BF> Claude, this is how you leave me a note and how I leave you one.
% BF> Anything after a % is invisible to LaTeX and never reaches the PDF.
% CL> I answer in place, under your note. Delete the pair when it's settled.

% --- The canonical figure block --------------------------------------------
% \label gives the figure a name YOU choose. \ref{fig:example} in the prose
% prints whatever number it ends up being. Insert a figure above this one and
% every number in the manuscript updates on the next build.
%
% Panel letters (A, B, ...) are drawn by make_figures.py, not by LaTeX --
% \subref{} exports to Word as an empty "()". See AGENTS.md.
\begin{figure}[tb]
  \centering
  \includegraphics[width=\linewidth]{fig1_example.pdf}
  \caption{\textbf{Short bold title of the figure.}
    \textbf{(A)} What panel A shows.
    \textbf{(B)} What panel B shows. Points are means of three replicates;
    error bars are one standard deviation.}
  \label{fig:example}
\end{figure}

As shown in Figure~\ref{fig:example}, \todo{the finding}. Note the tilde in
\verb|Figure~\ref{...}|: it is a non-breaking space, so "Figure" and its number
never get split across a line break.

% --- Citations --------------------------------------------------------------
% \citep{key} -> (Author et al. 2024).  \citet{key} -> Author et al. (2024).
% Keys come from refs.bib. Get a .bib entry from Google Scholar's "Cite"
% link, or from a DOI at doi2bib.org.
This is consistent with previous work \citep{example2024}.

% --- Equations --------------------------------------------------------------
% Numbered automatically in the PDF. But do NOT write \ref{eq:...} in the
% prose: pandoc numbers figures and tables, not equations, so it exports as
% the raw label "Equation [eq:model]". Refer to equations in words instead.
\begin{equation}
  y = \alpha + \beta x + \varepsilon
  \label{eq:model}
\end{equation}

\noindent
Parameters were estimated from the model above.

% --- Tables -----------------------------------------------------------------
% booktabs rules (\toprule/\midrule/\bottomrule) instead of vertical lines:
% this is what published tables look like. Column spec `l r r` = one left-
% aligned column, two right-aligned.
\begin{table}[tb]
  \centering
  \begin{threeparttable}
    \caption{Short description of the table.}
    \label{tab:example}
    \small
    \begin{tabular}{l r r}
      \toprule
      Group & Estimate & $P$ \\
      \midrule
      Treated & 0.42 & 0.001 \\
      Control & 0.08 & 0.640\tnote{a} \\
      \bottomrule
    \end{tabular}
    \begin{tablenotes}[flushleft]\footnotesize
      \item[a] Footnote markers attach to a cell and renumber themselves.
    \end{tablenotes}
  \end{threeparttable}
\end{table}

Table~\ref{tab:example} summarises \todo{what}.
